\chapter{Conclusiones} \label{sec:Conclusiones}

Los resultados experimentales arrojan las siguientes conclusiones por objetivo específico:

\begin{itemize}
    \item \textbf{OE1 (Rendimiento):} Calico entregó el mayor throughput TCP (1804.40 Mbps) y la menor latencia (12.54 ms). Flannel registró la menor tasa de retransmisión ---212 por GB transferido frente a 1\,021 de Calico, equivalentes a pérdidas del $0{,}029\%$ y el $0{,}145\%$---, si bien ambos valores se mantienen en el rango que TCP absorbe sin degradación perceptible. Ningún plugin resultó superior en todos los criterios: el liderazgo cambia según la métrica considerada, y esa ausencia de dominancia es el resultado que sostiene la necesidad de un modelo de decisión.
    
    \item \textbf{OE2 (Seguridad Zero Trust):} Calico, Cilium y Antrea aplican correctamente las Network Policies. Flannel presenta una vulnerabilidad arquitectónica al aceptar las políticas en la API de Kubernetes pero omitir su enforcement en el plano de datos, permitiendo tráfico no autorizado. Cilium fue la única solución capaz de aplicar reglas a nivel de capa 7 (HTTP). Este objetivo se alcanza de forma parcial: la evidencia demuestra el \emph{enforcement} efectivo de las políticas y acota su overhead, pero no cuantifica la reducción de la superficie de ataque como métrica numérica ni la probabilidad de incidentes de seguridad, según la delimitación de \cref{sec:delimitacion_oe2}.
    
    \item \textbf{OE3 (Modelo MCDA):} El algoritmo de normalización min-max y filtrado IQR resolvió las contrapartidas de rendimiento. Los tres perfiles evaluados arrojaron tres ganadores distintos: Cilium para Fintech (3.521) por su aplicación efectiva de políticas de capa 7, Calico para Streaming (3.567) por su throughput y latencia, y Flannel para IoT (3.259) por su menor huella de CPU/RAM. La ventaja de Flannel en IoT es estrecha frente a Cilium (3.148) y debe interpretarse como preferencia condicionada al ahorro de recursos, no como superioridad establecida. Se advierte, además, que el despliegue IoT en dispositivos ultra-restringidos demandará verificar la memoria aislada en reposo, al verse la métrica afectada por el ruido del host bajo carga.
    
    \item \textbf{OE4 (Prototipo Recomendador):} La evaluación multicriterio se integró de forma exitosa en una SPA. A partir de un cuestionario guiado de cinco preguntas, la herramienta recalcula dinámicamente el ranking sobre los datos procesados del testbed ---sin pérdida de precisión respecto de las mediciones crudas--- y entrega, para el CNI recomendado, el comando de instalación del plugin y una política base de denegación por defecto. La configuración avanzada de cada plugin y la actualización de datos en sesión quedan fuera de lo implementado.
\end{itemize}

La selección de un plugin CNI impacta directamente en el consumo de recursos de los nodos, la seguridad perimetral de los pods y el rendimiento de red. Una elección sustentada en datos, combinada con automatización GitOps, minimiza el sobredimensionamiento de infraestructura y el riesgo de brechas laterales.

\textbf{Impacto organizacional.} El beneficio para las organizaciones que adopten esta guía es cuantificable. La brecha de eficiencia medida entre el CNI más liviano y el más demandante ($9.6\%$ de CPU y $496.29$ MB de RAM por nodo) se traduce, a escala de un clúster de 50 nodos, en un ahorro aproximado de $9.6$ vCPU y $24.8$ GB de RAM ---capacidad que deja de adquirirse de forma preventiva y reduce el costo de infraestructura y energía. A su vez, el bajo overhead de la micro-segmentación Zero Trust (con una pérdida máxima de throughput del $4.35\%$ y variaciones de latencia acotadas) demuestra que una postura de seguridad verificable no exige sacrificar la calidad de servicio. Finalmente, al entregarse como un método replicable ---testbed como código, \emph{benchmarks} automatizados, modelo MCDA y un recomendador interactivo---, el proyecto permite que equipos sin especialización profunda en redes de Kubernetes reproduzcan esta evaluación en su propio entorno y sustituyan la intuición por evidencia telemática medida sobre el plano de datos.
